% sec-07-ai-peer-review.tex - Section 7, AI Peer Review.  A main section.
% FULL stage.  Figures 21, 22, 23 and 24 drawn; Tables 21 to 24 populated.
% Primary source is the AI peer review archive under new-trial-system/inputs;
% the role assignment comes from funding/RFA-RM-27-001-v2.
\section{AI Peer Review}
\label{sec:peerreview}

\subsection{The study the method comes from}

The AI peer review method used throughout this repository was established in a
study deposited November 30, 2025 \cite{aiprgliomatc}. That study was about
clinical trial patient matching in glioblastoma, not pancreatic cancer. The
disease is not what transfers; the review method is, and it transfers without
modification to the pancreatic cancer work reported in \S3 through \S6. The
study states its own central claim in one sentence: ``A primary benefit of AI
peer review is the ability to appropriately conduct iterative analyses and
corrections throughout the entire manuscript process. This responsible method
improves on over-worked and fatigued human recommendations provided after
manuscript completion'' \cite{aiprgliomatc}. That is the whole argument in
miniature, and everything else in this section is its measurement.

The prior system's timings are the study's own: peer review from the 1970s
through the 1990s meant mailing photographs of figures and waiting for
type-written comments; the late 1990s and 2000s brought electronic mail and web
portals ``to reduce physical mail delays by up to weeks, for a best-case review
cycle of 7-8 weeks''; and an average journal publication today ``can still take
several months for processing, including 1-2 rounds of review'', with faster
online journals at about one month \cite{aiprgliomatc}. Figure 21 lays those
timings against the new system's on one clock, and Table 21 gives the ratio on
each axis.

% ---------------------------------------------------------------------------
% FIGURE 21.  mermaid-type, sequenceDiagram.  Two lanes carrying six rows each
% over an identical 3.10 cm span, separated by one charcoal hairline, so the
% only difference the figure can express is the axis unit.
% ---------------------------------------------------------------------------
\begin{appfloat}[!tb]
\begin{appfig}[mermaid-type / sequenceDiagram]
\node[mmactor,text width=25mm] (au) at (0.9,0) {Author};
\node[mmactor,text width=25mm] (ed) at (5.0,0) {Editor, prior system};
\node[mmactor,text width=25mm] (hr) at (9.2,0) {Human reviewers, 2 to 3};
\node[mmactor,text width=25mm] (ai) at (13.5,0) {Model reviewers, 3 makers};
\foreach \x in {0.9,5.0,9.2,13.5}{\draw[mmlife] (\x,-0.55) -- (\x,-9.55);}
\begin{scope}[on background layer]
\fill[trialmistl] (-0.5,-4.95) rectangle (14.3,-0.75);
\end{scope}
\node[mmgray,text width=52mm,minimum height=4mm] at (5.05,-0.95) {Prior system: best case 7 to 8 weeks, typical several months};
\seqrow{-1.55}{0.9}{5.0}{mmmsg}{1 submit completed manuscript}
\seqrow{-2.15}{5.0}{9.2}{mmmsg}{2 invite reviewers, wait for acceptance}
\node[mmact,minimum height=1.20cm,fill=trialmist] at (9.2,-2.75) {};
\seqrow{-3.35}{9.2}{5.0}{mmret}{3 reports after weeks, sometimes months}
\seqrow{-3.95}{5.0}{0.9}{mmret}{4 decision plus reports, round 1}
\seqrow{-4.35}{0.9}{5.0}{mmmsg}{5 revision, round 2 begins}
\seqrow{-4.75}{5.0}{0.9}{mmret}{6 second decision, work already finished}
\draw[trialchar,line width=0.5pt] (-0.5,-5.10) -- (14.3,-5.10);
\node[mmsoft,text width=58mm,minimum height=4mm] at (7.20,-5.45) {New system: review during development, hour scale};
\seqrow{-6.05}{0.9}{13.5}{mmmsg}{7 draft plus dataset plus code, mid project}
\seqrow{-6.65}{13.5}{0.9}{mmret}{8 three independent reports, same day}
\seqrow{-7.25}{0.9}{13.5}{mmmsg}{9 corrections applied, re review requested}
\seqrow{-7.85}{13.5}{0.9}{mmret}{10 consensus and recorded disagreement}
\seqrow{-8.55}{0.9}{13.5}{mmmsg}{11 final pass before deposit}
\seqrow{-9.15}{13.5}{0.9}{mmret}{12 release recommendation, human decides}
\node[mmact,minimum height=3.10cm] at (13.5,-7.60) {};
\draw[decorate,decoration={brace,amplitude=4pt},trialchar,line width=0.5pt]
  (-0.75,-1.55) -- (-0.75,-4.75) node[midway,left=3pt,font=\tiny\sffamily,align=right] {weeks to\\months};
\draw[decorate,decoration={brace,amplitude=4pt},trialchar,line width=0.5pt]
  (-0.75,-6.05) -- (-0.75,-9.15) node[midway,left=3pt,font=\tiny\sffamily,align=right] {hours};
\pnote{-0.75}{-10.05}{Both lanes carry six rows over an identical 3.10 cm span. The row counts are
equal by construction, so the only\\thing the two braces can differ on is the
unit on the axis, which is the figure's whole argument.}
\end{appfig}
\vspace{-0.6cm}
\figcaption{Figure 21. One manuscript on two clocks: the prior system's serial
human\\rounds over months, and the new system's parallel model rounds within
hours.}
\end{appfloat}

\begin{apptable}
\begin{tabularx}{\textwidth}{>{\raggedright\arraybackslash}p{3.85cm} >{\raggedright\arraybackslash}p{3.55cm} >{\raggedright\arraybackslash}p{3.35cm} >{\raggedright\arraybackslash}X}
\headrow \thead{Economic axis} & \thead{Prior system} & \thead{New system} & \thead{Ratio}\\
\midrule
Latency, one round & 7 to 8 weeks best case & Same day & About 50 to 1 \\
Typical journal processing & Several months, 1 to 2 rounds & Hours, rounds unlimited & About 300 to 1 \\
Reviewers per round & 2 to 3 humans & 3 model manufacturers & 1 to 1 by count \\
Entry point & After completion & During development & Not comparable \\
Marginal cost per round & Reviewer time, unpriced & Inference cost, tens of dollars & Orders of magnitude \\
Artifacts absorbed per year & 2 to 6 per author & Over 30 deposited in 2026 & About 6 to 1 \\
Corrections before release & None, review follows release & Every round, before deposit & The decisive axis \\
\end{tabularx}
\end{apptable}
\vspace{-0.6cm}
\tabcap{Table 21. Six economic axes of one review round under both regimes,\\
and the seventh axis on which the two are not comparable at all.}

Figure 22 carries the same six axes as a grid, with the ratio computed in each
cell and the decisive seventh row separated by its own rule. A regime whose
review arrives after release cannot apply a correction to the released object,
whatever its cost or its latency, and that is why the seventh row is not an
economic axis at all.

% ---------------------------------------------------------------------------
% FIGURE 22.  d2-type, grid.  One row height, four fixed column widths, no
% edges, cut twice: once below the header and once above the decisive row.
% ---------------------------------------------------------------------------
\begin{appfloat}[!tb]
\begin{appfig}[d2-type / grid]
\def\rh{0.64}
\node[d2cellh,minimum width=4.4cm,minimum height=\rh cm,anchor=north west] at (0,0) {Economic axis};
\node[d2cellh,minimum width=4.0cm,minimum height=\rh cm,anchor=north west] at (4.4,0) {Prior system, human review};
\node[d2cellh,minimum width=4.0cm,minimum height=\rh cm,anchor=north west] at (8.4,0) {New system, AI review};
\node[d2cellh,minimum width=2.2cm,minimum height=\rh cm,anchor=north west] at (12.4,0) {Ratio};
\foreach \i/\a/\p/\n/\r/\rs in {%
  1/{Latency, one round}/{7 to 8 weeks best case}/{Same day}/{about 50 to 1}/{d2cellk},
  2/{Typical journal processing}/{Several months, 1 to 2 rounds}/{Hours, rounds unlimited}/{about 300 to 1}/{d2cellk},
  3/{Reviewers per round}/{2 to 3 humans}/{3 model manufacturers}/{1 to 1 by count}/{d2cellg},
  4/{Entry point}/{After completion}/{During development}/{not comparable}/{d2cellg},
  5/{Marginal cost per round}/{Reviewer time, unpriced}/{Inference cost, tens of dollars}/{orders of magnitude}/{d2cellk},
  6/{Artifacts absorbed per year}/{2 to 6 per author}/{over 30 deposited in 2026}/{about 6 to 1}/{d2cellk}}{
  \node[d2cellg,minimum width=4.4cm,minimum height=\rh cm,anchor=north west] at (0,{-\rh*\i}) {\a};
  \node[d2cell,minimum width=4.0cm,minimum height=\rh cm,anchor=north west] at (4.4,{-\rh*\i}) {\p};
  \node[d2celll,minimum width=4.0cm,minimum height=\rh cm,anchor=north west] at (8.4,{-\rh*\i}) {\n};
  \node[\rs,minimum width=2.2cm,minimum height=\rh cm,anchor=north west] at (12.4,{-\rh*\i}) {\r};}
\node[d2cellk,minimum width=4.4cm,minimum height=\rh cm,anchor=north west] at (0,{-\rh*7}) {Corrections before release};
\node[d2cellg,minimum width=4.0cm,minimum height=\rh cm,anchor=north west] at (4.4,{-\rh*7}) {None, review follows release};
\node[d2celll,minimum width=4.0cm,minimum height=\rh cm,anchor=north west] at (8.4,{-\rh*7}) {Every round, before deposit};
\node[d2cellh,minimum width=2.2cm,minimum height=\rh cm,anchor=north west] at (12.4,{-\rh*7}) {decisive};
\draw[trialchar,line width=0.7pt] (0,{-\rh}) -- (14.6,{-\rh});
\draw[trialchar,line width=0.7pt] (0,{-\rh*7}) -- (14.6,{-\rh*7});
\node[d2mid,text width=56mm,anchor=north west] at (0,{-\rh*8-0.18}) {Over 30 artifacts deposited in 2026 by one author};
\node[d2soft,text width=48mm,anchor=north west] at (6.20,{-\rh*8-0.18}) {Three manufacturers, one frozen artifact, one round};
\pnote{0}{-5.95}{Row 7 is separated by its own rule because it is not an economic axis: it is the
reason the economics matter. A\\regime whose review follows release cannot
correct the object that was released.}
\end{appfig}
\vspace{-0.6cm}
\figcaption{Figure 22. Six economic axes of one review round under both regimes,
with\\the ratio in each cell and the volume the prior system would have to
absorb.}
\end{appfloat}

\subsection{Three manufacturers, one human gate}

The pancreatic cancer funding application makes the role assignment explicit and
binding. Its provenance section states that the ``planned tripartisan workflow
assigns Claude Code the primary production role for code, regulatory
adaptations, protocols, IND materials, diagrams, and repository-scale artifacts.
ChatGPT Codex serves as the first independent peer reviewer, emphasizing
verification, validation, uncertainty quantification, external standards,
trial-literature synthesis, deep research, and PDF review. Google Gemini serves
as the second independent reviewer, emphasizing accelerated code problem
solving, meta-verification of prior stages, and administrative workflow
support'' \cite{rfarm27001v2}. The reviews ``occur at defined mid-project
milestones and after material corrections so that findings can change the work
before release, with the human PI retaining final authority''. And the limit is
stated in the same paragraph: these systems ``are disclosed as drafting and
peer-review aids; they are not applicants, investigators, sponsors, regulators,
clinicians, or decision-makers'' \cite{rfarm27001v2}. Table 22 names the three
manufacturers, the role each holds, and the human principal investigator who
retains final authority over all three, and Figure 23 draws where the
concurrency sits and where it stops: three manufacturers read one frozen
artifact and its recorded hash, so no branch can observe another's output and a
disagreement is evidence rather than an artifact of ordering.

\begin{apptable}
\begin{tabularx}{\textwidth}{>{\raggedright\arraybackslash}p{2.55cm} >{\raggedright\arraybackslash}p{3.05cm} >{\raggedright\arraybackslash}X}
\headrow \thead{Manufacturer} & \thead{Role} & \thead{What that role covers in this work}\\
\midrule
Anthropic & Primary production & Code, regulatory adaptations, protocols, IND materials, diagrams, repository-scale artifacts \\
OpenAI & First independent reviewer & Verification, validation, uncertainty quantification, external standards, trial-literature synthesis, deep research, PDF review \\
Google & Second independent reviewer & Accelerated code problem solving, meta-verification of prior stages, administrative workflow support \\
Human PI & Final authority & Source verification, regulatory accuracy, scientific interpretation, participant protection, final submission \\
\end{tabularx}
\end{apptable}
\vspace{-0.6cm}
\tabcap{Table 22. The three model manufacturers and the human principal\\
investigator, and the role each of them holds in the reviewed work.}

% ---------------------------------------------------------------------------
% FIGURE 23.  plantuml-type, activity with fork and join.  Three branches at a
% 5.20 cm pitch, 18 mm of clear canvas either side of each, one human gate.
% ---------------------------------------------------------------------------
\begin{appfloat}[!tb]
\begin{appfig}[plantuml-type / activity with fork and join]
\node[umlinit] (i0) at (7.50,0) {};
\node[umlbox,text width=46mm] (a1) at (7.50,-0.85) {Artifact reaches a defined mid project milestone};
\node[umlbox,text width=46mm] (a2) at (7.50,-2.05) {Freeze the artifact and its inputs, record the hash};
\node[umlbar,minimum width=12.8cm] (fk) at (7.50,-2.95) {};
\node[umlkey,text width=34mm] (b1) at (2.30,-4.45) {Anthropic, production role: code, regulatory text, diagrams};
\node[umlctrl,text width=34mm] (b2) at (7.50,-4.45) {OpenAI, first reviewer: verification, validation, uncertainty, standards};
\node[umlctrl,text width=34mm] (b3) at (12.70,-4.45) {Google, second reviewer: code problem solving and meta verification};
\node[umlbar,minimum width=12.8cm] (jn) at (7.50,-6.05) {};
\node[umlbox,text width=42mm] (a3) at (7.50,-6.85) {Join, three reports over one frozen artifact};
\node[mmdec,aspect=2.2,text width=24mm] (d1) at (7.50,-7.95) {Reports agree?};
\node[umlbox,text width=36mm] (cn) at (3.60,-9.05) {Apply the consensus corrections};
\node[umlstategray,text width=36mm] (dg) at (11.40,-9.05) {Record the disagreement verbatim, both positions};
\node[umlstategray,text width=36mm] (rt) at (11.40,-10.05) {Route to the Figure 24 resolution tree};
\node[umlkey,text width=46mm,line width=1pt] (pi) at (7.50,-11.15) {Human principal investigator reviews, approves or rejects};
\node[mmdec,aspect=2.2,text width=22mm] (d2) at (7.50,-12.35) {Approved?};
\node[umlbox,text width=32mm] (rl) at (4.20,-13.45) {Release, deposit, provenance recorded};
\node[umlfinal] (f1) at (4.20,-14.35) {};
\fill[trialburg] (f1.center) circle (1.1mm);
\node[umlbox,text width=32mm] (rj) at (11.40,-13.45) {Return to production with the rejection reason};
\node[umlfinal] (f2) at (11.40,-14.35) {};
\draw[trialchar,line width=0.6pt] ([shift={(-0.13,-0.13)}]f2.center) -- ([shift={(0.13,0.13)}]f2.center);
\draw[umlarrow] (i0) -- (a1);
\draw[umlarrow] (a1) -- (a2);
\draw[umlarrow] (a2) -- (fk);
\foreach \n in {b1,b2,b3}{\draw[umlarrow] (\n |- fk.south) -- (\n.north);
  \draw[umlarrow] (\n.south) -- (\n |- jn.north);}
\draw[umlarrow] (jn) -- (a3);
\draw[umlarrow] (a3) -- (d1);
\draw[umlarrow] (d1.west) -- ++(-0.55,0) -- (3.60,-7.95) -- node[umlguard,pos=0.6] {yes} (cn.north);
\draw[umlarrow] (d1.east) -- ++(0.55,0) -- (11.40,-7.95) -- node[umlguard,pos=0.6] {no} (dg.north);
\draw[umlarrow] (dg) -- (rt);
\draw[umlarrow] (cn.south) -- (3.60,-10.55) -- (pi.north west);
\draw[umlarrow] (rt.south) -- (11.40,-10.55) -- (pi.north east);
\draw[umlarrow] (pi) -- (d2);
\draw[umlarrow] (d2.west) -- ++(-0.55,0) -- (4.20,-12.35) -- node[umlguard,pos=0.6] {yes} (rl.north);
\draw[umlarrow] (d2.east) -- ++(0.55,0) -- (11.40,-12.35) -- node[umlguard,pos=0.6] {no} (rj.north);
\draw[umlarrow] (rl) -- (f1);
\draw[umldash] (rj) -- (f2);
\pnote{0}{-15.15}{The three branches sit at a 5.20 cm pitch against a 34 mm text width, leaving
18 mm of clear canvas either side, so\\the fork and join bars read as single
horizontal strokes rather than as a hatched band.}
\end{appfig}
\vspace{-0.6cm}
\figcaption{Figure 23. One artifact reviewed concurrently by three model
manufacturers,\\joined into one recorded disagreement set, and cleared only by the
human PI.}
\end{appfloat}

\subsection{What the review actually found}

The glioblastoma study is worth reporting in its numbers, because the numbers
are what make the method transferable rather than aspirational. A ten-file
dataset and a clinical BERT notebook were generated first, and after repeated
fixes reached ``a final notebook at a low test set performance of 67.3\%''. A
triple review conducted by three models from three manufacturers ``all yielded
the primary recommendation to upgrade to a tabular model, which was more suitable
for the csv dataset''; one of them then identified an open-source TabPFN-v2
classifier, and the recommendations were implemented to yield ``a test set
accuracy improvement to 94.0\% with few-shot learning along with peer review
corrections'' \cite{aiprgliomatc}. An evaluation standard generated for the
purpose scored the workflow at 87.5 percent overall quality and 0.831 efficiency
on human, prompt, time and cost metrics. Eleven attempts to produce a working
FastAPI repository failed before a two-model prompting strategy succeeded with
three correct patient trial match predictions. Three independent regulatory
analysts scored the six generated FDA, ICH and ISO documents at 8.0, 9.1 and 9.1
out of 10, and the study records the reason each gave rather than averaging them
\cite{aiprgliomatc}. Table 23 collects the measures.

\begin{apptable}
\begin{tabularx}{\textwidth}{>{\raggedright\arraybackslash}p{5.35cm} >{\raggedright\arraybackslash}p{2.55cm} >{\raggedright\arraybackslash}X}
\headrow \thead{Measure} & \thead{Value} & \thead{What changed it}\\
\midrule
Test-set accuracy before review & 67.3 percent & Initial clinical BERT pipeline \\
Test-set accuracy after review & 94.0 percent & Tabular model plus few-shot learning \\
Overall workflow quality & 87.5 percent & Evaluation standard, eight criteria \\
Efficiency score & 0.831 & Human, prompt, time and cost metrics \\
Regulatory analyst scores & 8.0, 9.1, 9.1 of 10 & Three manufacturers, recorded separately \\
Failed API attempts before success & 11 & Resolved by a two-model prompting strategy \\
Study duration, one author & 14 days & Dataset through FastAPI repository \\
\end{tabularx}
\end{apptable}
\vspace{-0.6cm}
\tabcap{Table 23. What the AI peer review changed, measured, in the glioblastoma\\study
whose method the pancreatic cancer work in this paper inherits.}

Figure 24 is the reason those three regulatory scores are reported separately
rather than as a mean of 8.7. None of the five terminal dispositions averages
the reports, and only one, record and proceed, leaves a disagreement without an
action; two of the five stop the artifact from progressing at all.

% ---------------------------------------------------------------------------
% FIGURE 24.  graphviz-type, decision tree.  Four ranks at a uniform 2.05 cm
% separation, deliberately asymmetric, with a 1.50 cm empty corridor between the
% two rank-3 subtrees.
% ---------------------------------------------------------------------------
\begin{appfloat}[!tb]
\begin{appfig}[graphviz-type / decision tree]
\node[gvboxk,text width=52mm,line width=0.9pt] (rt) at (7.50,0) {Three reports over one frozen artifact};
\node[mmdec,aspect=2.0,text width=22mm] (d1) at (7.50,-2.05) {Do all three agree?};
\node[gvboxs,text width=30mm] (t1) at (1.85,-4.10) {Apply the consensus};
\node[mmdec,aspect=2.0,text width=24mm] (d2) at (9.35,-4.10) {Factual or judgmental?};
\node[mmdec,aspect=2.0,text width=24mm] (d3) at (6.05,-6.15) {Checkable against a cited source?};
\node[mmdec,aspect=2.0,text width=24mm] (d4) at (12.65,-6.15) {Changes a safety claim?};
\node[gvboxs,text width=30mm] (t2) at (3.55,-8.20) {Resolve to the source, record which model erred};
\node[gvboxm,text width=30mm] (t5) at (7.35,-8.20) {Hold the artifact, re-run the round};
\node[gvboxm,text width=30mm] (t3) at (11.15,-8.20) {Escalate to the human PI, both positions verbatim};
\node[gvboxs,text width=30mm] (t4) at (14.15,-8.20) {Record and proceed, note in the limitations};
\draw[gvedge] (rt) -- (d1);
\draw[gvedge] (d1) -- node[umlguard,pos=0.5] {yes} (t1);
\draw[gvedge] (d1) -- node[umlguard,pos=0.5] {no} (d2);
\draw[gvedge] (d2) -- node[umlguard,pos=0.5] {factual} (d3);
\draw[gvedge] (d2) -- node[umlguard,pos=0.5] {judgmental} (d4);
\draw[gvedge] (d3) -- node[umlguard,pos=0.5] {yes} (t2);
\draw[gvedgeb] (d3) -- node[umlguard,pos=0.5] {no} (t5);
\draw[gvedgeb] (d4) -- node[umlguard,pos=0.5] {yes} (t3);
\draw[gvedge] (d4) -- node[umlguard,pos=0.5] {no} (t4);
\node[gvboxg,text width=44mm] at (0.35,-6.15) {No disposition averages the reports. Two of five stop the artifact.};
\pnote{0}{-9.45}{The tree is deliberately asymmetric: the left branch terminates in one step
because agreement is the common case,\\and the right branch in three because
disagreement is the case the figure exists to document.}
\end{appfig}
\vspace{-0.6cm}
\figcaption{Figure 24. One disagreement from three model reports to five terminal\\
dispositions, of which none is an average and only one is silent.}
\end{appfloat}

\subsection{Why only AI reviewers can carry this load}

The study's conclusion is not that AI review is convenient. It is that the
volume, complexity and frequency of AI-generated research have already passed
what human review can absorb: this next evolution of peer review marks ``the
pivot away from slow and irresponsible use of human peer review for
LLM-generated data, and back to fast automation of patient health applications
at scale'' \cite{aiprgliomatc}. The same study records that trust for the
process ``was established over a series of human-AI works reflected in LLM code
and non-code benchmarking, applications in literature, and prior author works
using AI in review roles coordinated across judging, external validation,
cross-verification, meta-verification, and tests according to FDA/ICH guidance''
\cite{aiprgliomatc, e2eefficiency2025, gbmsynergy2025}. Table 24 lists five
instances of that process across the author's deposited works, and the specific
change each produced.

\begin{apptable}
\begin{tabularx}{\textwidth}{>{\raggedright\arraybackslash}p{4.95cm} >{\raggedright\arraybackslash}p{1.85cm} >{\raggedright\arraybackslash}p{2.55cm} >{\raggedright\arraybackslash}X}
\headrow \thead{Work} & \thead{Date} & \thead{Reviewers} & \thead{What the review changed}\\
\midrule
Glioblastoma patient matching ML & Nov 30, 2025 & 3 manufacturers & Model class, 67.3 to 94.0 percent \\
Code generation competition & Dec 22, 2025 & 16 models & Iterative learning on FDA adverse events \\
NIH application v1.0 & Jul 7, 2026 & 2 reviewers & Structure and mechanism fit \\
NIH application v2.0 & Jul 12, 2026 & 2 reviewers & Budget, milestones, review schedule \\
Capitalization plan & Aug 11, 2026 & 2 reviewers & Column widths, caption balance, figure centering \\
\end{tabularx}
\end{apptable}
\vspace{-0.6cm}
\tabcap{Table 24. Five instances of AI peer review across the author's\\
deposited works, and the specific change each review produced.}

\subsection{What AI peer review is not}

Three limits are worth stating precisely, because the case for the method is
weaker if they are left implicit. First, model review is not regulatory review.
Nothing in \S7 substitutes for an institutional review board, an FDA
information request, or a data and safety monitoring board, and the funding
application says so in its own words: the models ``are not applicants,
investigators, sponsors, regulators, clinicians, or decision-makers''
\cite{rfarm27001v2}. Second, model review is not clinical validation.
Improvements over time on coding, scientific-reasoning and multimodal benchmarks
support current frontier models as independent reviewers, ``but not as clinical
validators'' \cite{rfarm27001v2, aiprgliomatc}. Third, model review does not
establish correctness by agreement. Three models agreeing is three reports, not
a proof, which is why Figure 24's tree never averages and why two of its five
dispositions stop the artifact rather than passing it.

The measured evidence for the second limit is in the source study itself. Its
own evaluation standard was ``only based on the current project, making its
applicability based more on similarities to the second model workflow'', and its
efficiency metrics ``were partly influenced by the author's priorities''
\cite{aiprgliomatc}. A method that reports that about itself is more useful than
one that does not, and it is the reason Table 23 reports three regulatory scores
separately with the reason each analyst gave rather than a single figure of
merit.

\subsection{Why the load only grows}

The volume argument is the one that does not reverse. In 2024 the author
deposited eight works; in 2025, twelve; in 2026 through August, more than
twenty, including an IND, two clinical trial protocols, five bill versions, ten
funding applications, two NIH application versions and a capitalization plan
\cite{repo2026}. Each of those is longer and more cross-referenced than its
predecessors, because each reads the ones before it. A review regime that took
seven to eight weeks per round in its best case was already the binding
constraint at eight works a year; at thirty it is not a constraint but an
exclusion, and the work simply does not enter it.

That is the sense in which the two systems are incompatible rather than merely
different. It is not that human review is worse. It is that human review is
sized for a production regime that no longer exists on the author's side of the
transaction, and no amount of goodwill inside the prior regime changes its
throughput. The new regime's reviewers are the only ones capable of handling
studies that are simultaneously larger, more complex and more frequent, and the
glioblastoma study's closing position is exactly that: the pivot is ``away from
slow and irresponsible use of human peer review for LLM-generated data, and back
to fast automation of patient health applications at scale''
\cite{aiprgliomatc}.
